% =====================================================================
% 1. Methodology (MLLB, Maria Martinez): experimental design.
% DECISIONS D3, D4, D14, D17, D18.
% =====================================================================
\section{Methodology}\label{sec:methodology}

\subsection{One protocol for every model}

The project compares architectures, so everything that is not the architecture is fixed and shared: the
data pipeline and split (Section~\ref{sec:mllb-data}), the loss, the training loop, the metrics and the run
logs (Section~\ref{sec:mllb-training}). A new model only has to be a \texttt{torch.nn.Module} that maps a
batch of shape (batch, 12, 1000) to five logits and to be registered by name; the same command trains it,
evaluates it and writes its log. In this way a difference between two models comes from the models and
not from two slightly different pieces of training code.

\subsection{Design of each block}

Each model block followed the same steps, and each step was a written decision before it was code:
\begin{enumerate}[leftmargin=1.4em,itemsep=1pt]
  \item \textbf{Baseline and grid.} A baseline and a list of configurations, each changing one factor with
        respect to a stated reference configuration. For the recurrent block (decision D17): 12
        configurations (cell, depth, direction, readout, attention score, attention crossed with cell,
        direction and length, input sequence). For the self-attention and Transformer block (D18): 15
        configurations, five of them rebuilding an encoder block one component at a time.
  \item \textbf{Time estimate.} Before the long runs, an \texttt{estimate} command measured the time per
        batch and the peak GPU memory of every configuration on the laptop, to check that the grid fitted
        the 4\,GB GPU and the time available.
  \item \textbf{Smoke test.} \texttt{all --fast} ran the whole pipeline on a few hundred records and two
        epochs; its outputs end in \texttt{\_fast}, are ignored by git and are rejected by the script that
        builds the numbers of the reports.
  \item \textbf{Validation runs.} Every configuration with seeds 42, 43 and 44, validation fold only.
  \item \textbf{Final models.} The selection rule fixed in step 1 picks the final model(s); each is trained
        again with \texttt{--final} and the three seeds, and only then evaluated on the test fold.
  \item \textbf{Report.} One command computes the summary, the paired comparisons, the tables and the
        figures from the logs and the saved predictions.
\end{enumerate}
The selection rules were: the best mean validation macro AUROC among all configurations (final RNN and
final Transformer), and the best among the five attention configurations (final RNN + attention).

\subsection{Comparisons and seeds}

A comparison between two configurations is always paired: same seed, same validation records, same
patient-bootstrap resamples (Section~\ref{sec:mllb-training}). For each seed we obtain the interval of the
difference in macro AUROC, and the tables report the mean difference and in how many seeds the interval
excludes 0 in each direction. Three seeds were chosen as the smallest number that lets us see whether an
effect repeats; with one seed, a difference within the noise of training would look like a result.

\subsection{Compute budget}

All runs used Adam with learning rate $10^{-3}$, batch 64, at most 50 epochs and early stopping with patience
5 on the validation macro AUROC (decision D14). Table~\ref{tab:mllb-cost} gives the cost of the main
configurations of my blocks, measured in the logs. The recurrent networks over 1000 raw steps were by far
the slowest, as expected from a sequential computation over ten times more steps; the Transformer, whose
attention over 100 tokens is computed in parallel, was cheaper than the LSTM over 1000 steps even with four
blocks.

\begin{table}[htbp]
\centering
\caption[Cost of the main configurations of my blocks]{Cost of the main configurations of my blocks
(mean over the three seeds, from the run logs): trainable parameters, sequence length, seconds per epoch,
best epoch and minutes per run.}
\label{tab:mllb-cost}
\small
\begin{adjustbox}{max width=\textwidth}
\begin{tabular}{llrrrrr}
\toprule
Block & Configuration & Parameters & Steps & s / epoch & Best epoch & min / run \\
\midrule
RNN & LSTM baseline (frames) & \rnnParamsBase{} & \rnnStepsBase{} & \onedec{\rnnEpochTimeMeanBase} & \onedec{\rnnBestEpochMeanBase} & \onedec{\rnnTrainTimeMinMeanBase} \\
RNN & GRU & \rnnParamsGru{} & \rnnStepsGru{} & \onedec{\rnnEpochTimeMeanGru} & \onedec{\rnnBestEpochMeanGru} & \onedec{\rnnTrainTimeMinMeanGru} \\
RNN & Bidirectional LSTM & \rnnParamsBi{} & \rnnStepsBi{} & \onedec{\rnnEpochTimeMeanBi} & \onedec{\rnnBestEpochMeanBi} & \onedec{\rnnTrainTimeMinMeanBi} \\
RNN & Convolutional front end & \rnnParamsConv{} & \rnnStepsConv{} & \onedec{\rnnEpochTimeMeanConv} & \onedec{\rnnBestEpochMeanConv} & \onedec{\rnnTrainTimeMinMeanConv} \\
RNN & Raw samples & \rnnParamsRaw{} & \rnnStepsRaw{} & \onedec{\rnnEpochTimeMeanRaw} & \onedec{\rnnBestEpochMeanRaw} & \onedec{\rnnTrainTimeMinMeanRaw} \\
RNN & Raw samples + attention & \rnnParamsRawAttn{} & \rnnStepsRawAttn{} & \onedec{\rnnEpochTimeMeanRawAttn} & \onedec{\rnnBestEpochMeanRawAttn} & \onedec{\rnnTrainTimeMinMeanRawAttn} \\
\midrule
Transformer & One self-attention layer & \transformerParamsSaSingle{} & \transformerTokensSaSingle{} & \onedec{\transformerEpochTimeMeanSaSingle} & \onedec{\transformerBestEpochMeanSaSingle} & \onedec{\transformerTrainTimeMinMeanSaSingle} \\
Transformer & Baseline (2 blocks) & \transformerParamsBase{} & \transformerTokensBase{} & \onedec{\transformerEpochTimeMeanBase} & \onedec{\transformerBestEpochMeanBase} & \onedec{\transformerTrainTimeMinMeanBase} \\
Transformer & 4 blocks & \transformerParamsLayersFour{} & \transformerTokensLayersFour{} & \onedec{\transformerEpochTimeMeanLayersFour} & \onedec{\transformerBestEpochMeanLayersFour} & \onedec{\transformerTrainTimeMinMeanLayersFour} \\
Transformer & Patches of 5 samples & \transformerParamsPatchFive{} & \transformerTokensPatchFive{} & \onedec{\transformerEpochTimeMeanPatchFive} & \onedec{\transformerBestEpochMeanPatchFive} & \onedec{\transformerTrainTimeMinMeanPatchFive} \\
Transformer & No positional encoding (final) & \transformerParamsNoPos{} & \transformerTokensNoPos{} & \onedec{\transformerEpochTimeMeanNoPos} & \onedec{\transformerBestEpochMeanNoPos} & \onedec{\transformerTrainTimeMinMeanNoPos} \\
\bottomrule
\end{tabular}
\end{adjustbox}
\end{table}
